\begin{abstract}
字节对编码分词器的词表与有序合并规则直接受到训练语料频率驱动，公开发布后可能泄露网站级训练成员信息。针对这一问题，本文提出安全聚合差分隐私批量字节对编码方法 SA-DP-BPE。该方法以完整网站为隐私单位，在公共候选域上对站点频率向量进行确定性 L1 截断，并利用 Paillier 加法同态加密完成多站点统计聚合；聚合服务器在密文域加入双边几何噪声，解密服务器仅获得带噪聚合并返回兼容的批量合并规则。在诚实但好奇且相互不共谋的双服务器威胁模型下，Paillier 的语义安全保证训练期统计的计算机密性，站点级差分隐私则约束最终分词器发布。实验在固定版本的 C4 网站语料上评估压缩率、词表重叠、频率估计、合并相似度和朴素贝叶斯五类分词器成员推断攻击，并在 AG News 上检验下游效用。标准 Plain BPE、协议匹配无 DP 基线和主 SA-DP-BPE 的五类攻击平均原始 AUC 分别为 \PlainMeanAttackAUC、\HEMeanAttackAUC{} 和 \PrimaryMeanAttackAUC；逐攻击校正分数方向后，相应平均 AUC 为 \PlainAdaptedAUC、\ProtocolBaselineAdaptedAUC{} 和 \PrimaryAdaptedAUC。主 SA-DP-BPE 的 AG News Macro-F1 为 \PrimaryMacroFOne，与 Plain BPE 的 \PlainMacroFOne{} 基本一致。在 \MeasuredCryptoCells{} 个 Paillier 参数组合的重复实测中，总耗时主要随客户端数量与候选维度的乘积增长，缩放模型拟合优度达到 $R^2=\CryptoTotalModelRtwo$。上述结果表明，SA-DP-BPE 在维持分词与下游效用的同时，降低了当前攻击设置下的发布侧成员风险，并实现了训练统计的密文聚合。
\end{abstract}

\begin{keywords}
分词器; 成员推断; 站点级差分隐私; Paillier同态加密; 安全聚合
\end{keywords}

\begin{eabstract}
The vocabulary and ordered merge rules of a byte-pair encoding (BPE) tokenizer are driven directly by training-corpus frequencies, so a released tokenizer may reveal website-level training membership. This paper proposes Secure-Aggregated Differentially Private Batched BPE (SA-DP-BPE). Treating a complete website as the privacy unit, SA-DP-BPE applies deterministic L1 clipping to site-frequency vectors over a public candidate domain and uses Paillier additive homomorphic encryption for multi-site aggregation. The aggregation server adds two-sided geometric noise in the ciphertext domain, while the decryption server obtains only the noisy aggregate and returns a compatible batch of merge rules. Under an honest-but-curious and non-colluding two-server threat model, Paillier semantic security provides computational confidentiality for training statistics, and site-level differential privacy constrains the released tokenizer. Experiments on a fixed revision of C4 websites evaluate five tokenizer membership-inference attacks---compression rate, vocabulary overlap, frequency estimation, merge similarity, and naive Bayes---together with downstream utility on AG News. The mean raw AUC values are \PlainMeanAttackAUC{} for standard Plain BPE, \HEMeanAttackAUC{} for the protocol-matched non-DP baseline, and \PrimaryMeanAttackAUC{} for the primary SA-DP-BPE configuration. After adapting the score direction separately for each attack, the corresponding mean AUC values are \PlainAdaptedAUC, \ProtocolBaselineAdaptedAUC, and \PrimaryAdaptedAUC. SA-DP-BPE achieves an AG News Macro-F1 of \PrimaryMacroFOne, compared with \PlainMacroFOne{} for Plain BPE. Across \MeasuredCryptoCells{} repeatedly measured Paillier configurations, total runtime scales primarily with the product of client count and candidate dimension, and the fitted scaling model reaches $R^2=\CryptoTotalModelRtwo$. Taken together, these results show that SA-DP-BPE preserves tokenizer and downstream utility while reducing release-side membership risk under the evaluated attack setting and enabling encrypted aggregation of training statistics.
\end{eabstract}

\begin{ekeywords}
tokenizer; membership inference; site-level differential privacy; Paillier homomorphic encryption; secure aggregation
\end{ekeywords}
